\color{unidarkgreen}{\subsection{Фиксация положения выключатея}}
\color{black}

\begin{enumerate}[label=\arabic{section}.\arabic{subsection}.\arabic*,  align=left, labelsep=4pt, leftmargin=0pt, itemindent=57pt]

\item
Функциональный блок фиксации положения выключателя (ФПВ) предназначен для определения
состояния ШСВ на основании технологической информации:

\begin{itemize}
    \item состояния блок-контактов выключателя;
    \item положения разъединителей в цепи выключателя;
    \item информации о выводе в ремонт выключателя.
\end{itemize}

\item
Контроль положения разъединителей при формировании сигналов включенного или отключенного
положения выключателя вводится программным ключом \textbf{«Контр. Р»}.

\item
Также функция ФПВ пофазно контролирует состояние блок-контактов выключателя и в случае
расхождения полюсов формирует сигнал \textit{«Внутр. пуск. ЗНФ В»}.

\vspace{2mm}
\begin{figure}[H]
\centering
\includegraphics[width=0.6\textwidth,height=0.6\textheight,keepaspectratio]{\fbpath/661.3381 ФПВ/img/FPV.pdf}
\captionof{figure}{Функциональная схема логики ФПВ}\label{pict_fpv:FPV}
\end{figure}

\small
\begin{longtable}{|>{\centering\arraybackslash}m{5.3cm}|>{\centering\arraybackslash}m{3.3cm}|>{\centering\arraybackslash}m{4.2cm}|>{\centering\arraybackslash}m{1.8cm}|>{\centering\arraybackslash}m{1cm}|}
\caption{Перечень уставочных параметров ФПВ\hfill\vspace{-0.5\baselineskip}}\label{fpv:tbl1}\\ 
\hline
\rowcolor{gray!30}
Параметр & Условное обозначение на схеме & Значение/ Диапазон & Единица измерения & Шаг \\ 
\hline
\endfirsthead
\caption*{\hspace{3pt}\emph{Продолжение таблицы \ref{fpv:tbl1}\hfill\vspace{-0.5\baselineskip}}} \\
\hline
\rowcolor{gray!30}
Параметр & Условное обозначение на схеме & Значение/ Диапазон & Единица измерения & Шаг \\ 
%\hline
\endhead

%===m>ФПВ|-
%===m
\end{longtable}
\normalsize

\end{enumerate}